\chapter{Minimum Spanning Trees}

\begin{definition}
Let $G=(V,E)$ be a connected undirected graph with edge weights $w:E\to\R$. A \emph{spanning tree} is a subset $T\subseteq E$ forming a tree on all of $V$
($|T|=n-1$). A \emph{minimum spanning tree} (MST) minimizes $w(T)=\sum_{e\in T}w(e)$.
\end{definition}
Applications: network design (cheapest way to connect sites), clustering (delete the $k-1$ heaviest MST edges to get a max-spacing $k$-clustering),
and approximation algorithms (metric TSP, Chapter 11).

\section{The cut property}
A \emph{cut} is a partition $(S,V\setminus S)$ with $\emptyset\ne S\ne V$. An edge \emph{crosses} the cut if it has one endpoint on each side.
A set $A$ of edges \emph{respects} the cut if no edge of $A$ crosses it.

\begin{theorem}[Cut property]
Let $A$ be a subset of some MST, $(S,V\setminus S)$ a cut respecting $A$, and $e$ a minimum-weight edge crossing the cut. Then $A\cup\{e\}$ is contained in
some MST. (In particular, if edge weights are distinct, the lightest edge crossing any cut is in \emph{every} MST.)
\end{theorem}
\begin{proof}[Proof (exchange)]
Let $T$ be an MST containing $A$. If $e=uv\in T$ we are done. Otherwise $T\cup\{e\}$ has a unique cycle, which crosses the cut at $e$ and so must cross it
again at some edge $e'\in T$, $e'\notin A$ (since $A$ respects the cut). Then $T'=T\cup\{e\}\setminus\{e'\}$ is a spanning tree (removing an edge of the cycle
keeps it connected) with $w(T')=w(T)+w(e)-w(e')\le w(T)$ since $w(e)\le w(e')$. So $T'$ is an MST containing $A\cup\{e\}$.
\end{proof}

\begin{theorem}[Cycle property]
If $e$ is a strictly heaviest edge on some cycle, then $e$ is in no MST.
\end{theorem}

\begin{important}[Generic MST algorithm]
Start with $A=\emptyset$. While $|A|<n-1$, find a cut respecting $A$ and add a lightest crossing edge. By the cut property, $A$ is always contained in an MST.
Kruskal and Prim are two ways of choosing the cut.
\end{important}

\section{Kruskal's algorithm}
\begin{lstlisting}[style=pseudo]
Kruskal(G, w):
    A = {}
    for each v in V do MakeSet(v)
    sort E by weight, increasing
    for each edge (u, v) in sorted order do
        if Find(u) $\ne$ Find(v) then          // different components
            A = A $\cup$ {(u, v)}; Union(u, v)
    return A
\end{lstlisting}
\textbf{Correctness}: when $(u,v)$ is added, let $S$ be the component of $u$ in $(V,A)$. The cut $(S,V\setminus S)$ respects $A$, and $(u,v)$ is the lightest
crossing edge: any lighter crossing edge would have been considered earlier and added (its endpoints were in different components then, and still are).
Edges that are skipped would close a cycle, so the result is a forest, and since $G$ is connected it ends with $n-1$ edges.

\subsection{Union-find (disjoint sets)}
Operations: \code{MakeSet(x)}, \code{Find(x)} (return the representative of $x$'s set), \code{Union(x,y)}.
\begin{itemize}
  \item \textbf{Linked lists with union by size}: each element stores a pointer to its representative, so \code{Find} is $\bigO(1)$; \code{Union} relabels the smaller set.
        An element's set at least doubles every time it is relabeled, so each element is relabeled $\le\log_2 n$ times: total union cost $\bigO(n\log n)$.
  \item \textbf{Forest with union by rank and path compression}: $m$ operations take $\bigO(m\,\alpha(n))$, where $\alpha$ is the inverse Ackermann function ($\le4$ in practice).
\end{itemize}
\begin{lstlisting}[style=pseudo]
Find(x):
    if parent[x] $\ne$ x then parent[x] = Find(parent[x])    // path compression
    return parent[x]
Union(x, y):
    a = Find(x); b = Find(y)
    if rank[a] < rank[b] then swap a, b
    parent[b] = a
    if rank[a] == rank[b] then rank[a] = rank[a] + 1
\end{lstlisting}
\textbf{Kruskal runtime}: sorting $\bigO(m\log m)=\bigO(m\log n)$ (since $m\le n^2$), plus $2m$ finds and $n-1$ unions, which is $\bigO(m\log n)$ even with the simple list
structure. Total $\bigO(m\log n)$.

\section{Prim's algorithm}
Prim grows a single tree from a root $r$, always adding the lightest edge leaving the tree (cut $=(\text{tree},\text{rest})$).
\begin{lstlisting}[style=pseudo]
Prim(G, w, r):
    for each v do key[v] = $\infty$; parent[v] = nil
    key[r] = 0
    Q = min-priority queue of all vertices keyed by key
    while Q not empty do
        u = ExtractMin(Q)
        for each v adjacent to u do
            if v in Q and w(u, v) < key[v] then
                parent[v] = u; DecreaseKey(Q, v, w(u, v))
    return {(v, parent[v]) : v $\ne$ r}
\end{lstlisting}
$\mathrm{key}[v]$ is the weight of the lightest edge from $v$ to the current tree. Correctness is directly the cut property with $S$ = the extracted vertices.

\textbf{Runtime}: $n$ \code{ExtractMin} and up to $m$ \code{DecreaseKey}.
\begin{center}
\begin{tabular}{lll}
\toprule
Priority queue & Total time & Best for \\
\midrule
Unsorted array & $\bigO(n^2)$ & dense graphs \\
Binary heap & $\bigO((n+m)\log n)=\bigO(m\log n)$ & sparse graphs \\
Fibonacci heap & $\bigO(m+n\log n)$ & theory \\
\bottomrule
\end{tabular}
\end{center}

\begin{example}
Vertices $a,b,c,d$ with edges $ab{:}1$, $bc{:}2$, $ac{:}3$, $cd{:}4$, $bd{:}5$.
Kruskal: take $ab$ (1), $bc$ (2), skip $ac$ (cycle), take $cd$ (4): weight 7.
Prim from $a$: add $ab$, then lightest edge leaving $\{a,b\}$ is $bc$, then $cd$. Same tree.
\end{example}

\section{Uniqueness and other facts}
\begin{itemize}
  \item If all weights are distinct, the MST is unique.
  \item Adding a constant to all weights, or applying any increasing function, does not change which trees are MSTs (only the order of edges matters).
        This is \emph{not} true for shortest paths.
  \item Negative weights are fine for MST algorithms.
  \item Bor\r{u}vka's algorithm: in each round every component adds its lightest outgoing edge; $\bigO(\log n)$ rounds, $\bigO(m\log n)$ total.
\end{itemize}

\section*{Summary}
\begin{itemize}
  \item Cut property: the lightest edge across a cut respecting $A$ can be added safely; proved by exchanging with a cycle edge.
  \item Kruskal: sorted edges + union-find, $\bigO(m\log n)$.
  \item Prim: grow one tree using a priority queue, $\bigO(m\log n)$ with a binary heap, $\bigO(n^2)$ with an array.
\end{itemize}
